\documentclass[a4paper]{article}

%------------------------------------------------
% الحزم
%------------------------------------------------

\usepackage{enumitem}
\usepackage[margin=2.2cm]{geometry}
\usepackage{amsmath,amssymb}
\usepackage{graphicx}
\usepackage{array}
\usepackage{fancyhdr}
\usepackage[table]{xcolor}
\usepackage{tikz}
\usetikzlibrary{positioning,arrows.meta,shapes.geometric,calc}
\usepackage[skins]{tcolorbox}
\usepackage{polyglossia}

%------------------------------------------------
% اللغات والخطوط
%------------------------------------------------

\setdefaultlanguage[
    calendar=gregorian,
    numerals=maghrib
]{arabic}

\setotherlanguage{english}

% إذا لم تكن الخطوط مثبتة في الجهاز، تُؤخذ من المجلد fonts/
\IfFontExistsTF{Amiri}{%
    \newfontfamily\arabicfont[Script=Arabic, Scale=1.15]{Amiri}
}{%
    \newfontfamily\arabicfont[
        Script=Arabic, Scale=1.15,
        Path=fonts/, Extension=.ttf,
        UprightFont=*-Regular, BoldFont=*-Bold,
        ItalicFont=*-Italic, BoldItalicFont=*-BoldItalic
    ]{Amiri}
}

\IfFontExistsTF{AlBattar}{%
    \newfontfamily\arabicfontsf[Script=Arabic, Scale=1.2]{AlBattar}
    % خط العناوين (يعمل حتى بعد استعمال \textenglish)
    \newfontfamily\khatt[Script=Arabic, Scale=1.2]{AlBattar}
}{%
    \newfontfamily\arabicfontsf[
        Script=Arabic, Scale=1.2, Path=fonts/,
        BoldFont=ae_AlBattar.ttf
    ]{ae_AlBattar.ttf}
    % خط العناوين (يعمل حتى بعد استعمال \textenglish)
    \newfontfamily\khatt[
        Script=Arabic, Scale=1.2, Path=fonts/,
        BoldFont=ae_AlBattar.ttf
    ]{ae_AlBattar.ttf}
}

%------------------------------------------------
% الألوان: لون رئيسي هادئ + لون تنبيه فقط
%------------------------------------------------

\definecolor{asasi}{HTML}{3D6A8E}   % اللون الرئيسي: للرسوم والمخططات
\definecolor{tanbih}{HTML}{C0603A}  % لون التنبيه: للتحذيرات والحالات السلبية
\definecolor{ramadi}{HTML}{5F6B73}  % رمادي للنصوص الثانوية والأسهم
\definecolor{dakin}{HTML}{263238}   % لون العناوين

% أسماء قديمة تُحوَّل إلى اللونين فقط
\colorlet{azraq}{asasi}
\colorlet{fayruzi}{asasi}
\colorlet{banafsaji}{asasi}
\colorlet{akhdar}{asasi}
\colorlet{ahmar}{tanbih}
\colorlet{burtuqali}{asasi}

%------------------------------------------------
% الصناديق
%------------------------------------------------

\tcbset{
    asl/.style={
        enhanced,
        arc=2mm,
        boxrule=0.6pt,
        colback=white,
        colframe=black!18,
        colbacktitle=black!5,
        coltitle=dakin,
        titlerule=0pt,
        left=8pt, right=8pt, top=5pt, bottom=5pt,
        fonttitle=\khatt\bfseries,
        before skip=8pt, after skip=8pt
    }
}

% أهداف الدرس
\newtcolorbox{ahdaf}{asl, title={أهداف الدرس}}

% تعريف (خط رئيسي رفيع على الجانب)
\newtcolorbox{taarif}[1][تعريف]{asl, title={#1},
    borderline east={2pt}{0pt}{asasi}}

% مثال
\newtcolorbox{mithal}[1][مثال]{asl, title={#1}}

% تنبيه (اللون الوحيد البارز في النص)
\newtcolorbox{intabih}[1][انتبه!]{asl, colback=tanbih!5, colframe=tanbih,
    colbacktitle=tanbih, coltitle=white, title={#1}}

% ركن أستاذ اللغة العربية
\newtcolorbox{arabiyya}[1][في قسم اللغة العربية]{asl, title={#1}}

% نص أمر يكتب للأداة
\newtcolorbox{amr}[1][ramadi]{enhanced, arc=1mm, boxrule=0pt,
    colback=black!4, borderline east={2.5pt}{0pt}{#1},
    left=8pt, right=10pt, top=4pt, bottom=4pt,
    before skip=6pt, after skip=6pt}

% الخلاصة
\newtcolorbox{khulasa}{asl, title={خلاصة الدرس}}

% عنوان القسم
\newcounter{qism}
\newcommand{\qism}[1]{%
    \stepcounter{qism}%
    \par\bigskip
    {\large\khatt\bfseries{\color{asasi}\arabic{qism}.}\ {\color{dakin}#1}}\par
    \vspace{-2pt}{\color{black!20}\rule{\linewidth}{0.6pt}}\par\nopagebreak[4]
    \smallskip}

% عنوان فرعي
\newcommand{\far}[1]{\par\smallskip
    {\color{dakin}\khatt\bfseries {\color{asasi}$\blacktriangleleft$}\ #1}\par\nopagebreak[4]}

% كلمة مهمة (خط عريض فقط)
\newcommand{\kalima}[2][]{\textbf{#2}}

%------------------------------------------------
% أنماط الرسوم
%------------------------------------------------

\newcommand{\nodecolor}{}

\tikzset{
    % كتابة النص العربي من اليمين إلى اليسار داخل العقد
    % (لون النص يُعطى بالمفتاح lawn=... وليس text=...)
    every node/.append style={
        execute at begin node={\nodecolor\everypar{\beginR}\beginR},
        execute at end node=\endR},
    lawn/.code={\def\nodecolor{\color{#1}}},
    khatwa/.style={draw=#1, fill=#1!12, very thick, rounded corners=6pt,
        minimum height=1.2cm, text width=2.3cm, align=center,
        font=\khatt\bfseries},
    khatwa/.default=azraq,
    sahm/.style={-{Stealth[length=3mm]}, very thick, color=ramadi},
    raqm/.style={circle, fill=#1, lawn=white, inner sep=2pt,
        font=\small\bfseries},
}

%------------------------------------------------
% إعداد الصفحة
%------------------------------------------------

\renewcommand{\footrulewidth}{1pt}
\renewcommand{\headrulewidth}{1pt}

\rhead{\khatt المدرسة العليا للأساتذة ببوزريعة}
\lhead{\khatt ملحقة بومرداس}

\cfoot{\thepage}
\lfoot{\khatt الأستاذ: د. بن عبد الرزاق يوسف}

\parindent 0pt
\setlength{\headsep}{15pt}
\setlength{\parskip}{4pt}

\pagestyle{fancy}

\setlist[enumerate,1]{label=\arabic*), itemsep=2pt, topsep=3pt}
\setlist[itemize]{itemsep=2pt, topsep=3pt}

% الجداول
\arrayrulecolor{black!22}
\renewcommand{\arraystretch}{1.35}

% رأس جدول ملون
\newcommand{\ras}[1]{\cellcolor{black!7}\color{dakin}\bfseries #1}

% عنوان الدرس
\newcommand{\unwan}[3]{%
    \begin{center}
        {\color{asasi}\rule{\linewidth}{1.2pt}}\par\medskip
        {\khatt\color{ramadi} #1}\par\smallskip
        {\Large\khatt\bfseries\color{dakin} #2}\par\smallskip
        {\khatt\color{ramadi} #3}\par\medskip
        {\color{asasi}\rule{\linewidth}{1.2pt}}
    \end{center}}

%================================================
% بداية الوثيقة
%================================================

\begin{document}

%------------------------------------------------
% العنوان
%------------------------------------------------

\unwan{درس الفصل الثاني}
{كيف يتعلم الذكاء الاصطناعي؟}
{البيانات ــ التدريب ــ الاختبار}

\begin{ahdaf}
في نهاية هذا الدرس، تستطيع أن:
\begin{itemize}
    \item تشرح مراحل تعلم الآلة: من البيانات إلى النتيجة.
    \item تفرّق بين بيانات التدريب وبيانات الاختبار.
    \item تفهم لماذا يجب أن تكون البيانات متنوعة.
    \item تتعرف على الانحياز الخوارزمي وطرق التقليل منه.
\end{itemize}
\end{ahdaf}

%================================================
\qism{التلميذ والآلة: تشبيه بسيط}
%================================================

كيف يتعلم التلميذ الإعراب؟
يقرأ الأستاذ معه أمثلة كثيرة مُعرَبة، فيلاحظ التلميذ القاعدة.
ثم يعطيه الأستاذ جملًا جديدة ليتأكد أنه فهم فعلًا.

الآلة تتعلم بطريقة قريبة من هذه.
تأخذ أمثلة كثيرة، وتكتشف فيها \kalima{الأنماط}، أي الأشياء التي تتكرر.

\begin{center}
\begin{tikzpicture}[
    t/.style={draw=#1, fill=#1!10, thick, rounded corners=5pt,
        text width=5.6cm, align=center, minimum height=1cm}]
    \node[font=\khatt\bfseries, lawn=ramadi] at (3.4,1) {التلميذ};
    \node[font=\khatt\bfseries, lawn=azraq]     at (-3.4,1) {الآلة};
    \node[t=ramadi] (t1) at (3.4,0)    {يقرأ أمثلة مُعرَبة};
    \node[t=ramadi] (t2) at (3.4,-1.4) {يلاحظ القاعدة};
    \node[t=ramadi] (t3) at (3.4,-2.8) {يُعرب جملًا جديدة};
    \node[t=azraq] (a1) at (-3.4,0)    {تأخذ بيانات التدريب};
    \node[t=azraq] (a2) at (-3.4,-1.4) {تكتشف الأنماط};
    \node[t=azraq] (a3) at (-3.4,-2.8) {تُختبر ببيانات جديدة};
    \draw[sahm] (t1)--(t2); \draw[sahm] (t2)--(t3);
    \draw[sahm] (a1)--(a2); \draw[sahm] (a2)--(a3);
    \foreach \i in {1,2,3}
        \draw[dashed, thick, ramadi!60] (t\i) -- (a\i);
\end{tikzpicture}
\end{center}

%================================================
\qism{مراحل التعلم}
%================================================

يمر تعلم الآلة بخمس مراحل متتابعة:

\begin{center}
\begin{tikzpicture}[node distance=5mm]
    \node[khatwa=azraq, text width=2cm]      (a) {البيانات};
    \node[khatwa=fayruzi, text width=2cm, left=of a]   (b) {التدريب};
    \node[khatwa=banafsaji, text width=2cm, left=of b] (c) {النموذج};
    \node[khatwa=burtuqali, text width=2cm, left=of c] (d) {الاختبار};
    \node[khatwa=akhdar, text width=2cm, left=of d]    (e) {النتيجة};
    \draw[sahm] (a)--(b); \draw[sahm] (b)--(c);
    \draw[sahm] (c)--(d); \draw[sahm] (d)--(e);
\end{tikzpicture}
\end{center}

\begin{center}
\rowcolors{2}{black!3}{white}
\begin{tabular}{>{\centering\arraybackslash}m{2.6cm}
                >{\raggedleft\arraybackslash}m{11.5cm}}
\hline
\ras{المرحلة} & \ras{دورها} \\
\hline
\kalima{البيانات}      & أمثلة كثيرة: نصوص أو صور أو أصوات. \\
\kalima[fayruzi]{التدريب}  & يتعلم النظام من البيانات، ويكتشف العلاقات بينها. \\
\kalima[banafsaji]{النموذج} & هو ما ينتج عن التدريب: برنامج «تعلّم» القاعدة. \\
\kalima[burtuqali]{الاختبار} & نجرب النموذج على أمثلة جديدة لم يرها من قبل. \\
\kalima[akhdar]{النتيجة}   & توقع أو تصنيف أو إجابة يقدمها النموذج. \\
\hline
\end{tabular}
\end{center}

%================================================
\qism{بيانات التدريب وبيانات الاختبار}
%================================================

\begin{taarif}
\textbf{بيانات التدريب:} البيانات التي يتعلم منها النموذج ويكتشف فيها الأنماط.

\textbf{بيانات الاختبار:} بيانات جديدة نستعملها لتقييم النموذج بعد تدريبه.
\end{taarif}

نقسم البيانات عادة إلى جزأين: جزء كبير للتدريب، وجزء صغير للاختبار.

\begin{center}
\begin{tikzpicture}
    \fill[azraq!70, rounded corners=3pt] (0,0) rectangle (-11.2,1);
    \fill[asasi!45, rounded corners=3pt] (-11.2,0) rectangle (-14,1);
    \node[lawn=white, font=\khatt\bfseries] at (-5.6,0.5)
        {بيانات التدريب (80\%)};
    \node[lawn=white, font=\khatt\bfseries] at (-12.6,0.5)
        {الاختبار (20\%)};
\end{tikzpicture}
\end{center}

\begin{intabih}[لماذا لا نختبر بالأمثلة نفسها؟]
لو اختبر الأستاذ التلميذ بالجمل نفسها التي شرحها في القسم،
فقد يجيب التلميذ من الذاكرة لا من الفهم.
كذلك النموذج: نختبره ببيانات جديدة لنعرف هل \textbf{فهم} حقًا.
\end{intabih}

%================================================
\qism{تنوع البيانات}
%================================================

يتعلم النموذج ممّا يراه فقط.
إذا كانت البيانات ناقصة أو غير متوازنة، كان التعلم ناقصًا.

\begin{mithal}
لدينا 500 صورة لهواتف، و500 صورة لحواسيب، و100 صورة فقط لألواح رقمية.
\end{mithal}

\begin{center}
\begin{tikzpicture}[y=0.006cm]
    \draw[thick, ramadi] (0.5,0) -- (-8.5,0);
    \fill[azraq!70]     (0,0)  rectangle (-2,500);
    \fill[fayruzi!70]   (-3,0) rectangle (-5,500);
    \fill[ahmar!70]     (-6,0) rectangle (-8,100);
    \node[above, font=\bfseries] at (-1,500) {500};
    \node[above, font=\bfseries] at (-4,500) {500};
    \node[above, font=\bfseries] at (-7,100) {100};
    \node[below, font=\khatt] at (-1,0) {هواتف};
    \node[below, font=\khatt] at (-4,0) {حواسيب};
    \node[below, font=\khatt] at (-7,0) {ألواح رقمية};
\end{tikzpicture}
\end{center}

صور الألواح الرقمية قليلة جدًا.
لذلك سيخطئ النموذج كثيرًا عندما يرى لوحًا رقميًا جديدًا.
نقول إن هذه البيانات \kalima[ahmar]{غير ممثِّلة}.

\begin{taarif}[البيانات الجيدة]
البيانات الجيدة \textbf{متنوعة} و\textbf{ممثِّلة}:
تغطي مختلف الحالات التي سيواجهها النظام في الواقع.
\end{taarif}

%================================================
\qism{الانحياز الخوارزمي}
%================================================

\begin{taarif}
\textbf{الانحياز الخوارزمي} نتائج غير عادلة أو غير دقيقة،
سببها مشكلة في البيانات أو في طريقة تصميم النظام.
\end{taarif}

\begin{center}
\begin{tikzpicture}[node distance=9mm]
    \node[khatwa=ahmar, text width=3.2cm] (a) {بيانات ناقصة\\أو من فئة واحدة};
    \node[khatwa=burtuqali, text width=3.2cm, left=of a] (b) {نموذج تعلّم\\صورة ناقصة};
    \node[khatwa=banafsaji, text width=3.2cm, left=of b] (c) {نتائج منحازة\\وغير عادلة};
    \draw[sahm] (a)--(b); \draw[sahm] (b)--(c);
\end{tikzpicture}
\end{center}

\begin{arabiyya}[مثال من قسم اللغة العربية]
نموذج لتصحيح الأخطاء تدرّب على نصوص صحفية حديثة فقط.
عندما نعطيه بيتًا من الشعر الجاهلي، قد يعدّ بعض الكلمات الفصيحة القديمة أخطاءً.
السبب: لم يرَ هذا النوع من النصوص أثناء التدريب.
\end{arabiyya}

\far{كيف نقلل الانحياز؟}
\begin{itemize}
    \item نجمع بيانات \textbf{متنوعة} تمثل جميع الفئات والحالات.
    \item نوازن بين عدد الأمثلة في كل فئة.
    \item نختبر النموذج على حالات مختلفة، ونراجع أخطاءه.
\end{itemize}

\begin{intabih}
لا يكفي أن نقول: «النموذج أعطى النتيجة، إذن هي صحيحة».
النموذج يعكس البيانات التي تعلّم منها، بما فيها من نقص.
\end{intabih}

\begin{khulasa}
\begin{itemize}
    \item تتعلم الآلة من \textbf{البيانات}، كما يتعلم التلميذ من الأمثلة.
    \item المراحل: البيانات $\leftarrow$ التدريب $\leftarrow$ النموذج
          $\leftarrow$ الاختبار $\leftarrow$ النتيجة.
    \item نختبر النموذج دائمًا ببيانات \textbf{جديدة} لم يرها.
    \item البيانات الناقصة تؤدي إلى \textbf{الانحياز}، فنحتاج بيانات متنوعة وممثِّلة.
\end{itemize}
\end{khulasa}

\end{document}
